Realistyczne odwzorowanie zachowania cieczy od lat pozostaje jednym
z~najbardziej wymagających zadań grafiki komputerowej. Woda towarzyszy
człowiekowi na co dzień, dlatego widz natychmiast dostrzega każde uproszczenie:
nienaturalny rozprysk, brak załamania światła czy zbyt gładką powierzchnię fali.
Jednocześnie zapotrzebowanie na wiarygodne symulacje płynów stale rośnie ---
od produkcji filmowej i~gier komputerowych, przez wizualizacje naukowe,
po inżynierskie analizy przepływów. W~produkcji offline, gdzie pojedyncza klatka
może powstawać godzinami, jakość symulacji ogranicza wyłącznie budżet
obliczeniowy; w~zastosowaniach interaktywnych ograniczenie jest znacznie
bardziej rygorystyczne, ponieważ pełny krok fizyki wraz z~renderowaniem musi
zmieścić się w~kilkunastu milisekundach.

Spełnienie tego wymagania stało się możliwe dzięki ewolucji procesorów
graficznych, które z~wyspecjalizowanych układów rasteryzujących przekształciły
się w~masowo równoległe maszyny obliczeniowe ogólnego przeznaczenia.
Współczesne GPU wykonuje jednocześnie dziesiątki tysięcy wątków, co naturalnie
odpowiada strukturze symulacji cząsteczkowych, w~których ten sam zestaw operacji
stosowany jest niezależnie do każdego elementu cieczy. Pełne wykorzystanie tego
potencjału wymaga jednak narzędzi dających programiście bezpośrednią kontrolę
nad sprzętem: niskopoziomowych API graficznych oraz języków systemowych
łączących wydajność z~gwarancjami poprawności. Tymczasem dojrzałe systemy
symulacji cieczy --- zarówno pakiety animacyjne, jak i~biblioteki badawcze ---
wykonują centralne obliczenia w~przeważającej mierze na procesorze głównym,
a~obecność karty graficznej wykorzystują głównie do wyświetlania wyników.

Celem niniejszej pracy jest zaprojektowanie i~implementacja interaktywnego
silnika symulacji cieczy, w~którym kompletny potok obliczeniowy --- od dynamiki
cząstek, przez wymuszanie nieściśliwości, po wolumetryczną wizualizację ---
wykonywany jest w~czasie rzeczywistym w~całości na procesorze graficznym.
Fundament fizyczny systemu stanowi cząsteczkowa metoda SPH w~wariancie DFSPH,
wymuszającym nieściśliwość cieczy przez dwuetapową korekcję gęstości
i~dywergencji pola prędkości. Warstwę technologiczną tworzą język programowania
Rust oraz niskopoziomowe API graficzne Vulkan --- połączenie oferujące
statyczne bezpieczeństwo pamięci przy zachowaniu jawnej kontroli nad zasobami
sprzętowymi. Zakres pracy obejmuje sformułowanie podstaw teoretycznych metody
SPH i~algorytmu DFSPH, uzasadniony dobór stosu technologicznego, implementację
solvera oraz potoku wizualizacji wraz z~autorskimi rozwiązaniami
architektonicznymi, a~także ilościową analizę wydajności i~zbieżności
zaimplementowanych algorytmów.

Praca składa się z~trzech rozdziałów. W~rozdziale pierwszym przedstawiono
podstawy teoretyczne: porównanie eulerowskiego i~lagrange'owskiego opisu ruchu
płynów, równania Naviera-Stokesa, dyskretyzację jądrową metody SPH wraz
z~algorytmem DFSPH oraz koncepcję wolumetrycznej wizualizacji metodą przemarszu
promienia. Rozdział drugi zawiera przegląd istniejących rozwiązań symulacji
cieczy oraz charakterystykę wybranych technologii: języka Rust, API Vulkan
i~biblioteki Vulkano, narzędzi diagnostycznych, a~ponadto analizę porównawczą
funkcji wygładzających, struktur wyszukiwania sąsiedztwa i~modeli zjawisk
fizycznych. Rozdział trzeci dokumentuje implementację systemu: architekturę
aplikacji, organizację danych w~pamięci GPU, potok wyszukiwania sąsiadów wraz
z~porównaniem dwóch algorytmów sortowania, kolejne etapy solvera DFSPH oraz
potok renderowania --- każdorazowo z~uzasadnieniem podjętych decyzji
i~pomiarami wydajności. Pracę zamyka podsumowanie zestawiające osiągnięte
rezultaty oraz kierunki dalszego rozwoju systemu.
